%---------------------------------------------------------------------
\chapter{Introduction}
\label{ch:intro}

\section{Background}
Post-training with reinforcement learning has shifted from human-preference reward
models (RLHF~\cite{ouyang2022instructgpt}) toward \emph{verifiable} rewards, where
a programmatic checker --- an exact-match grader for GSM8K, a unit test for code
--- supplies a binary or scalar reward with no learned reward model in the
loop~\cite{lambert2024tulu3,guo2025r1}. Within this RLVR setting, Group Relative
Policy Optimization (GRPO)~\cite{shao2024deepseekmath} has become the default
algorithm. It drops PPO's value network~\cite{schulman2017ppo} and instead
estimates the advantage of each sampled completion \emph{relative to its own group}
of same-prompt samples. Because GRPO removes a whole model (the critic), it is
cheaper and simpler to run at scale, and it was central to the DeepSeek-R1
reasoning models~\cite{guo2025r1}; it is now the recipe most open post-training
efforts reach for first.

This simplicity has a cost that is rarely reported. Because the advantage is a
within-group standardisation, any prompt for which every sampled completion
receives the \emph{same} reward contributes an advantage of exactly zero, and
therefore \emph{no policy-gradient (advantage) signal} --- only the KL penalty
toward the reference policy still acts. On an easy dataset the model quickly solves most
prompts, so a large fraction of every batch collapses to zero advantage. We call
this the \textbf{Zero-Variance Fraction (ZVF)}, and it is, empirically, the single
largest measured source of zero-advantage sampling in this fixed study. ZVF is the namesake and the
organising measurement of this project: every study in the portfolio either
measures it, tests a lever intended to reduce it, or uses it as a feature.

\section{A Motivating Example}
Consider a single training step with a group size of $G=8$ on an easy GSM8K prompt
that the current policy already solves. All eight sampled completions produce the
correct final answer, so each receives reward $r_i=1$. The group mean is $\mu_g=1$
and the group standard deviation is $\sigma_g=0$, so by Eq.~\eqref{eq:adv} every
advantage $A_i=0$. The advantage (surrogate) term of the objective then contributes
no reward-derived surrogate gradient for the entire group of eight rollouts --- eight full forward-decode
passes that cost real tokens and money --- so the group supplies \emph{no
task-learning signal}. (The KL-regularisation term of Eq.~\eqref{eq:loss} still nudges
the policy toward the reference; the point is that no \emph{reward-driven} learning
happens.)
Now imagine a batch of $16$ such prompts where $12$ are already solved: $12/16=0.75$
of the batch is ``dark'' compute. This is not a corner case; Chapter~\ref{ch:results}
shows it is the \emph{typical} regime ($\zvf\approx0.72$--$0.77$). The Zero-Variance
Fraction turns this intuition into one loggable number, and the rest of the report is
an attempt to measure it precisely and to test whether it can be reduced into real
held-out accuracy.

\section{Problem Statement}
Three concrete problems motivate this project.
\begin{enumerate}[leftmargin=1.6em,itemsep=3pt]
  \item \textbf{A reproducibility crisis.} GRPO/RLVR results are reported on
  different frameworks with different (often undocumented) reward parsers, decoding
  temperatures, KL coefficients, and clip ranges. A headline number cannot be
  attributed to the \emph{algorithm} versus the \emph{stack}.
  \item \textbf{Advantage-signal starvation is invisible.} Standard training logs
  report reward and loss, but not what fraction of the batch supplied a non-zero
  reward-relative advantage. Such rollouts may still receive KL regularisation, but
  they provide no task-reward direction and still consume sampling compute.
  \item \textbf{No provenance standard.} There is no machine-readable record of
  \emph{what} was trained, on which verifier version, with which rollouts, so
  contamination and silent reward drift go undetected.
\end{enumerate}

\section{Objectives}
The Phase-1 objectives, against which the work is evaluated, are:
\begin{enumerate}[leftmargin=1.6em,itemsep=2pt]
  \item Build a \textbf{unified harness} that runs the identical task, reward, and
  decoding configuration across multiple RL post-training back-ends, so that
  cross-framework differences are attributable to the stack.
  \item Instrument a \textbf{per-step telemetry layer} that surfaces ZVF, gradient
  utilisation, entropy, completion length, and KL, and that saves raw group reward
  tensors so any published diagnostic can be \emph{recomputed}, not merely trusted.
  \item \textbf{Measure} the zero-variance phenomenon on real reward tensors, and
  test whether any re-baselining or curriculum lever can convert zero-advantage
  sampling into held-out accuracy.
  \item Establish a \textbf{scientific protocol} --- multi-seed, matched baselines,
  held-out evaluation with reported $n$, deterministic recomputation --- and apply it
  honestly, reporting negative and null results where they occur.
  \item Scope the most \emph{original} Phase-2 direction (a GRPO provenance /
  reporting standard) via a literature scan for open white space.
\end{enumerate}

\section{Scope and Non-Goals}
Phase-1 is deliberately a \emph{measurement and diagnostics} phase on small-to-mid
open models (Qwen2.5 / Qwen3.5 at $1.5$B--$4$B) and math/code verifiers (GSM8K,
a HumanEval subset). It is explicitly \textbf{not} an attempt to top a GRPO
leaderboard, nor to train a frontier reasoning model. Held-out evaluation sets are
intentionally small and every gain is reported with its denominator so that the
reader can see when a result is noise-limited. Novel-algorithm claims (P4 length
weighting, P7 controller) are \emph{designed} in Phase-1 and deferred to Phase-2
for powered evaluation. This honest scoping is itself a deliverable: it is what
makes the benchmark trustworthy.

\section{Contribution of the Candidate}
\label{sec:contribution}
This is an \emph{individual Semester-4 continuation}. The candidate's new work in
this report comprises the ZVF-centred diagnostic framing, post-capstone experiment
expansion, telemetry and recomputation analyses, the P1--P8 synthesis, provenance
prototype, and the writing and evaluation presented here. The scientific posture is
\emph{verify $\rightarrow$ replicate where feasible $\rightarrow$ report the honest
null}.

\paragraph{Inherited foundation vs.\ new work.} The repository intentionally retains
the Semester-3 Group-6 foundation: the multi-framework research codebase, literature
survey, baseline integrations, and early experiments are inherited shared
infrastructure frozen at Git tag \texttt{capstone-final-2026-04-25}. Semester 4 begins
after that boundary and is the candidate's solo continuation under the present guide.
This report does not claim the inherited Group-6 foundation as new individual work;
it claims the post-boundary diagnostic, experiment, analysis, and provenance additions
that build on it. \texttt{PROJECT\_\-HISTORY.md} and \texttt{sem 4 work/\allowbreak PROVENANCE.md}
record the boundary.

\paragraph{What was built vs.\ leveraged.} To delineate the candidate's own effort from
reused infrastructure: the ZVF / collapse instrumentation, the
recompute-from-tensors diagnostics, the process-per-run campaign launcher, the P8
method-trace classifier, and the P5 provenance prototype are post-boundary additions
used by this project. Leveraged (not claimed as new individual work) are the shared
Semester-3 codebase and baseline experiments, the underlying model weights
(Qwen2.5 / Qwen3.5), the managed training back-end (Tinker), the RL framework libraries
being unified (TRL, veRL, OpenRLHF), and standard libraries (PyTorch, Transformers,
PEFT, scikit-learn). Table~\ref{tab:impl} summarises the implementation status.

\begin{table}[h]
\centering\small
\begin{tabular}{@{}L{0.34\linewidth}L{0.12\linewidth}L{0.44\linewidth}@{}}
\toprule
\textbf{Component} & \textbf{Status} & \textbf{Evidence} \\
\midrule
GRPO reference loop (\texttt{platform\_\-tinker/\allowbreak tinker\_\-rl/\allowbreak grpo.py}) & Done & P2/P3/campaign runs \\
ZVF / collapse instrumentation & Done & recomputed by \texttt{p2\_\-collapse\_\-analysis.py} \\
Per-step telemetry $\to$ JSONL / W\&B & Done & project \texttt{rlvr-openings} \\
Process-per-run campaign launcher & Done & 9-run multi-seed campaign \\
White-box P1 layer profiler & Done & 1.5B + scaled 3B re-test \\
P8 method-trace classifier & Prototype & \texttt{detector\_\-reproduced.json} (row-CV AUROC $0.838$) \\
Per-framework adapters (TRL/veRL/OpenRLHF) & Partial & implemented; four-way run is Phase-2 \\
P5 provenance / P6 registry & Prototype & schema + \texttt{minreport verify} \\
P4 length weighting, P7 controller & Designed & specified for Phase-2 \\
\bottomrule
\end{tabular}
\caption{Implementation-completion matrix. The core measurement and diagnostic stack is
complete; cross-framework comparison and the designed algorithmic levers (P4, P7) are
Phase-2 work. Overall the implemented core corresponds to the ${\sim}80\%$
Phase-1 code expectation, with the remaining ${\sim}20\%$ scoped for Phase-2.}
\label{tab:impl}
\end{table}

\section{Project Plan and Timeline}
The project follows the PES Phase-1 review cadence: a zeroth review (title, abstract,
literature, proposed system, timeline), three internal reviews at roughly fortnightly
intervals, and an End-Semester Assessment (ESA) viva. Table~\ref{tab:timeline} maps the
deliverables of this report onto that schedule.

\begin{table}[h]
\centering\small
\begin{tabular}{@{}L{0.20\linewidth}L{0.24\linewidth}L{0.46\linewidth}@{}}
\toprule
\textbf{Milestone} & \textbf{Window (2026)} & \textbf{Deliverable} \\
\midrule
Initiation / Zeroth & May & Title, abstract, literature survey, proposed system \\
First Review & early June & Architecture, algorithm design, ${\sim}30\%$ code \\
Second Review & mid June--July & Detailed design, intermediate results, ${\sim}80\%$ code \\
Third Review & early July & Overall design, experimental results, ${\sim}100\%$ code, paper draft \\
Phase-1 ESA & July & Presentation + code walkthrough/demo (viva) \\
Phase-2 & Jul--Oct & Novelty (P5 provenance), publication \\
\bottomrule
\end{tabular}
\caption{Phase-1 milestones and deliverables against the PES review schedule.}
\label{tab:timeline}
\end{table}

\section{Organisation of the Report}
Chapter~\ref{ch:lit} surveys the relevant literature and positions the eight
studies against the closest prior art. Chapter~\ref{ch:srs} specifies the system
requirements (functional, non-functional, test cases, and the hardware/software
environment). Chapter~\ref{ch:method} describes the proposed methodology --- the
unified harness architecture, the telemetry design, the formalisation of ZVF, and
the P1--P8 portfolio design. Chapter~\ref{ch:impl} details the implementation.
Chapter~\ref{ch:results} presents the intermediate results and discussion,
including the rigour and threats-to-validity analysis. Chapter~\ref{ch:conclusion}
concludes and lays out the Phase-2 plan. Appendices give the complete run registry
and reproducibility details.

%---------------------------------------------------------------------
